\subsection{Conditional traces and core minorants}
\label{app:core-conditional-traces}

Conditional foreign closures at an independent base have an exact
description by a minorant. They need not themselves have a common
constant-free realization.

Let $I$ be a finite set of at least two foreign owners. Write
$Q=\{0<\alpha<\beta\}$ for the algebra $3_m$ of
\cref{app:core-dissemilattice}, with $+$ the chain maximum and
$uv=0$ if either operand is zero, and $uv=\max(u,v)$ otherwise.
At coordinate $s$, interpret $\star_s$ as multiplication and every
other choice as $+$. Let $P_I$ be the product of these coordinate
algebras, with named coordinatewise union; $P_{I,0}$ additionally
names its all-zero tuple. This zero is a trace value, not an empty
winner-menu constant. Named union is available in the full core with
one additional owner $\rho$.

For loopless closures $d_s$ on a finite nonempty set $B$, let $f_d$
be their coordinatewise trace function. At coordinate $s$ put
$J_s(C)=B\setminus d_s(B\setminus C)$. On inputs zero outside $C$,
$\beta$ on $M\subseteq C$, and $\alpha$ on $C\setminus M$, its value
is zero if $J_s(C)$ is empty, $\beta$ if $J_s(C)\cap M\ne\varnothing$,
and $\alpha$ otherwise. Each coordinate has the one-owner expression
\[
 f_{d,s}=\sum_{\substack{H\ne\varnothing\\B\setminus H\text{ closed}}}
                 \prod_{j\in H}x_j.
\]
It therefore has a $P_{I,0}$ term: lift this expression using
$\star_s$, attach zero at each foreign owner $t\ne s$ to isolate
coordinate $s$, and sum the isolated coordinates.

\begin{theorem}[Conditional trace characterization]
\label{core:thm:conditional-trace-minorants}
The following conditions on $(d_s)_{s\in I}$ are equivalent:
\begin{enumerate}
\item There is a core operation over $I\cup\{\rho\}$, an independent
  base $B$ for $c_\rho$, and $D=c_\rho(B)\setminus B$, such that
  $d_s(E)=c_s(E\cup D)\cap B$ for every $s\in I$ and $E\subseteq B$.
\item There is a constant-free foreign-choice/union word $g$ using
  every variable of $B$, with $g\leq f_d$ pointwise on $P_I$ in the
  product chain order.
\end{enumerate}
In the reverse construction one extra variable suffices: $D$ can
be a singleton.
\end{theorem}
The trace inequality is equivalently $c_s^g\supseteq d_s$ for every
foreign owner. Extending $d$ by the identity at $\rho$, the source-factor
formula also gives $g\leq F_d$ on winner menus.
\begin{proof}
Take the neutral residual on $U=c_\rho(B)$, whose closures restrict
to $c_s(E)\cap U$ by \cref{core:thm:neutral-elimination}.
Its $\rho$-profile is positive at $B$ and negative at every
$B\setminus\{j\}$: independence excludes $j$ from that closure.
Prune each $\rho$-owned node to a child positive at $B$; every foreign
node and union node keeps both children. The remaining leaves lie in
$B$, and the resulting foreign-choice/union word $g$ uses all of $B$.
Otherwise its profile would be positive at some $B\setminus\{j\}$,
where pruning can only decrease the original profile.

Pruning decreases the $\rho$-player view at every input. Indeed the
$\rho$-owned operation is join in that view, and each foreign gate
is monotone in its order: foreign atoms lie below $\nv<\pv$.
If $j\in d_s(E)\setminus E$, use the residual's pure-$s$ probe with
$j=\{\rho\}$, $E\cup D$ at $\{s\}$, and everything else at the
foreign menu $I$. Its output omits $\rho$.
The pruned output does too, and $g$ has no $D$ inputs. Hence
$c_s^g(E)\supseteq d_s(E)$. Its trace interior is therefore contained
in $J_s$, which is exactly the pointwise trace bound $g\leq f_d$.

Conversely choose a $P_{I,0}$ term $F$ for $f_d$, using the displayed
isolators, and replace it by $F+0$. Introduce a new variable $Y$ in
place of every zero, giving a constant-free word $T_0$, and put
\[
 T=T_0\star_\rho g.
\]
Each old choice and named union is join in the $\rho$-trace. Thus,
writing $b=\sum_{j\in B}x_j$, the $\rho$-trace of $T$ is
$(b+Y)b$. Its interior on $C\subseteq B\cup\{Y\}$ is $C$ when
$C\cap B\ne\varnothing$, and empty otherwise. Consequently
\[
 c_\rho(E)=
 \begin{cases}E\cup\{Y\},&B\subseteq E,\\E,&\text{otherwise}.
 \end{cases}
\]
The base $B$ is independent and $c_\rho(B)\setminus B=\{Y\}$.
At an old coordinate $s$, set $Y$ to trace zero. The root is join
there, giving $f_{d,s}+g_s=f_{d,s}$, exactly the required conditional
closure. Interpreting named unions by their core term gives an ordinary
constant-free realization. Their $\rho$-trace remains join, including
when that union term uses $\rho$ internally.
\end{proof}

\begin{proposition}[Conditional data need not be jointly realizable]
\label{core:prop:conditional-trace-nonrealizable}
An actual four-owner core operation can have, at an independent base,
three conditional foreign closures with no common core realizer.
\end{proposition}
\begin{proof}
Use foreign owners $0,1,2$ and $B=\{x,y,z,w\}$, with the rules
\[
 d_0:\{x,z,w\}\Rightarrow y,\qquad
 d_1:\{y,z\}\Rightarrow x,\qquad d_2:\{x,y\}\Rightarrow z.
\]
These loopless closures have no common realizer. At the independent
$0$-base $\{x,z,w\}$ their conditional foreign closures, with
$D=\{y\}$, contain $z\Rightarrow x$ and $x\Rightarrow z$.
This violates \eqref{core:eq:conditional-mutual-capture}, a necessary
consequence of (FE).

Their coordinate trace formulas are
\[
 f_0=m+ym,\quad m=x+z+w;\qquad
 f_1=w+n+xn,\quad n=y+z;\qquad
 f_2=w+p+zp,\quad p=x+y.
\]
There is nevertheless a full-support minorant. Put
\[
 h=y+(x\star_1(y\star_2z))+(z\star_2(y\star_1x)),
 \qquad g=w+(w\star_0h).
\]
At coordinates $1,2$, distribution in $Q$ gives $g_s=f_s$.
At coordinate $0$, $h=x+y+z$, so $g_0=0$ when $w=0$ and is the
full input join otherwise. The function $f_0$ is zero only when
$x=z=w=0$ and is the full join otherwise. Hence $g\leq f_d$.
The preceding theorem, using the isolated-coordinate term $F$ and
replacing its zeros by $Y$, supplies a five-variable core word
$T_0\star_\rho g$ with precisely these conditional closures.
\end{proof}
